% ============================================================
%  PAUTA DE RESPUESTAS
%  Finanzas Públicas del Perú
%  Fuentes: MMM 2026-2029 (Cuadros C-1, C-6, C-10) y BCRP (C-4)
%           The fiscal history of Peru in two graphs (MEF / BCRP)
% ============================================================

\documentclass[12pt, a4paper]{article}

% -- Codificación e idioma ------------------------------------
\usepackage[utf8]{inputenc}
\usepackage[T1]{fontenc}
% En Overleaf (texlive completo) descomentar:
\usepackage[spanish, es-tabla]{babel}

% -- Márgenes ------------------------------------------------
\usepackage[top=2.5cm, bottom=2.5cm, left=3cm, right=3cm]{geometry}

% -- Tipografía ----------------------------------------------
\usepackage{lmodern}
\usepackage{microtype}

% -- Matemáticas ---------------------------------------------
\usepackage{amsmath, amssymb}

% -- Tablas --------------------------------------------------
\usepackage{booktabs}
\usepackage{multirow}
\usepackage{array}
\usepackage[table]{xcolor}
\definecolor{headgray}{gray}{0.88}
\definecolor{rowgray}{gray}{0.95}

% -- Listas --------------------------------------------------
\usepackage{enumitem}

% -- Hipervínculos --------------------------------------------
\usepackage[colorlinks=true, linkcolor=blue!60!black,
            urlcolor=blue!70!black, citecolor=gray]{hyperref}

% -- Encabezados y pies --------------------------------------
\usepackage{fancyhdr}
\setlength{\headheight}{14.5pt}
\pagestyle{fancy}
\fancyhf{}
\rhead{\small Finanzas Públicas del Perú — Pauta de Respuestas}
\lhead{\small \thepage}
\renewcommand{\headrulewidth}{0.4pt}

% -- Interlineado --------------------------------------------
\usepackage{setspace}
\setstretch{1.15}

% -- Colores de énfasis --------------------------------------
\newcommand{\emp}[1]{\textbf{#1}}
\newcommand{\nota}[1]{\textit{\small #1}}

% -- Cajas de interpretación ---------------------------------
\usepackage{tcolorbox}
\tcbset{colback=blue!5!white, colframe=blue!45!black,
        fonttitle=\bfseries, boxrule=0.6pt, arc=3pt}

% -- Título --------------------------------------------------
\title{\Large\bfseries Pauta de Respuestas\\[6pt]
       \large Finanzas Públicas del Perú\\[4pt]
       \normalsize Marco Macroeconómico Multianual 2026--2029 \&
       Serie Histórica BCRP}
\author{}
\date{\small Mayo 2026}

% ============================================================
\begin{document}
\maketitle
\thispagestyle{fancy}

\tableofcontents
\vspace{1cm}

% ============================================================
\section{¿Cómo están nuestras finanzas públicas?
         \texorpdfstring{\\[-4pt]}{}\normalsize
         Respuesta a los incisos (c), (d) y (e)}
% ============================================================

% ------------------------------------------------------------
\subsection{La aritmética del déficit \texorpdfstring{---}{---}
            inciso (c)}
% ------------------------------------------------------------

\subsubsection*{Los tres indicadores y su identidad}

El resultado económico, el resultado primario y los intereses del Sector Público
No Financiero (SPNF) no son indicadores independientes: están unidos por una
identidad contable que conviene tener siempre presente.
Definamos cada concepto:

\begin{itemize}[leftmargin=1.5em]
    \item \emp{Resultado primario (RP):} diferencia entre los ingresos del
          gobierno general y sus gastos no financieros (es decir, excluye el
          pago de intereses de la deuda pública). Mide el esfuerzo fiscal
          \emph{propio} del período, independientemente de la herencia de deuda.
    \item \emp{Intereses ($i$):} servicio de la deuda pública acumulada.
          Representan una obligación rígida derivada de decisiones de
          endeudamiento pasadas, no de la política fiscal corriente.
    \item \emp{Resultado económico (RE):} balance global del sector público.
          Es la medida más completa de la posición fiscal en un año dado.
\end{itemize}

La identidad que los vincula es:

\begin{equation}
    \boxed{\mathrm{RE} = \mathrm{RP} - i}
    \label{eq:identidad}
\end{equation}

\noindent donde todas las variables se expresan como porcentaje del PBI.
La ecuación~\eqref{eq:identidad} implica que un gobierno puede tener superávit
primario y aun así registrar un déficit económico si el servicio de deuda
heredado supera dicho superávit.

\subsubsection*{Verificación con los datos del MMM 2026--2029}

El Cuadro~1 del repositorio permite verificar la identidad para cada año del
horizonte de proyección:

\begin{table}[h!]
\centering
\caption{Resultado primario, intereses y resultado económico del SPNF
         (\% del PBI), 2024--2029}
\label{tab:aritmetica}
\setlength{\tabcolsep}{9pt}
\renewcommand{\arraystretch}{1.25}
\begin{tabular}{lrrrl}
\toprule
\rowcolor{headgray}
\textbf{Año} &
\textbf{Res.\ primario} &
\textbf{Intereses} &
\textbf{Res.\ económico} &
\textbf{Verificación} \\
 & \textit{(RP)} & \textit{(i)} & \textit{(RE = RP$-$i)} & \\
\midrule
2024 & $-1{,}78$ & $1{,}67$ & $-3{,}46$ & $-1{,}78 - 1{,}67 = -3{,}45\;\checkmark$ \\
\rowcolor{rowgray}
2025 & $-0{,}49$ & $1{,}66$ & $-2{,}15$ & $-0{,}49 - 1{,}66 = -2{,}15\;\checkmark$ \\
2026 & $-0{,}07$ & $1{,}72$ & $-1{,}80$ & $-0{,}07 - 1{,}72 = -1{,}79\;\checkmark$ \\
\rowcolor{rowgray}
2027 & $+0{,}36$ & $1{,}76$ & $-1{,}40$ & $+0{,}36 - 1{,}76 = -1{,}40\;\checkmark$ \\
2028 & $+0{,}78$ & $1{,}78$ & $-1{,}00$ & $+0{,}78 - 1{,}78 = -1{,}00\;\checkmark$ \\
\rowcolor{rowgray}
2029 & $+0{,}75$ & $1{,}74$ & $-0{,}99$ & $+0{,}75 - 1{,}74 = -0{,}99\;\checkmark$ \\
\bottomrule
\end{tabular}
\smallskip

\nota{Fuente: elaboración propia sobre la base del Cuadro~1 del MMM 2026--2029.
Las pequeñas diferencias en la verificación se deben a redondeo.}
\end{table}

\subsubsection*{¿Cuándo el resultado primario deja de ser negativo?}

Según las proyecciones del MMM, el resultado primario del SPNF pasa de terreno
negativo a positivo en \emp{2027}, cuando registra $+0{,}36\%$ del PBI.
Ese es el primer año en que el gobierno, con sus ingresos corrientes, cubre
íntegramente sus gastos no financieros.

\begin{tcolorbox}[title={Interpretación clave}]
Que el resultado primario sea positivo en 2027 pero el resultado económico siga
siendo \emph{negativo} ($-1{,}40\%$ del PBI) no es una contradicción: significa
que el gobierno ya no necesita endeudarse para financiar su gasto corriente y de
capital, pero el stock de deuda acumulado en años anteriores genera un flujo
de intereses ($1{,}76\%$ del PBI en 2027) que aún supera ese superávit primario.
En otras palabras, \emp{el déficit en 2027 es enteramente producto de la herencia
de deuda pasada}, no de una nueva indisciplina fiscal.
\end{tcolorbox}

% ------------------------------------------------------------
\subsection{Resultado convencional vs.\ resultado estructural ---
            inciso (d)}
% ------------------------------------------------------------

\subsubsection*{La diferencia conceptual}

El resultado económico \emp{convencional} registra los ingresos y gastos
\emph{efectivamente observados} en el año, que incluyen tanto sus componentes
permanentes como sus componentes cíclicos (por ejemplo, la mayor recaudación
del IR cuando los precios de las materias primas están elevados, o los ingresos
extraordinarios de ventas de activos).

El resultado económico \emp{estructural}, en cambio, ajusta los ingresos para
neutralizar dos efectos:

\begin{enumerate}[label=(\roman*), leftmargin=2.5em]
    \item \emp{Brecha del producto:} cuando el PBI observado supera el PBI
          potencial (brecha positiva), la recaudación tributaria es
          transitoriamente alta porque las empresas obtienen mayores ganancias
          y el consumo es elevado. El ajuste elimina ese exceso.
    \item \emp{Ingresos extraordinarios:} se excluyen ingresos no recurrentes
          como grandes ventas de activos de no domiciliados, según la
          metodología aprobada por la Resolución Ministerial
          N.\textordmasculine~024-2016-EF/15.
\end{enumerate}

El resultado estructural mide, por tanto, la posición fiscal que tendría el
gobierno si la economía operara en su nivel potencial. Es el indicador más
adecuado para evaluar el \emp{sesgo de la política fiscal discrecional} con
independencia del ciclo económico.

\subsubsection*{La brecha entre ambos indicadores y lo que revela}

\begin{table}[h!]
\centering
\caption{Resultado económico convencional vs.\ estructural del SPNF
         (\% del PBI / PBI potencial), 2024--2029}
\label{tab:estructural}
\setlength{\tabcolsep}{9pt}
\renewcommand{\arraystretch}{1.25}
\begin{tabular}{lrrr}
\toprule
\rowcolor{headgray}
\textbf{Año} &
\textbf{RE convencional} &
\textbf{RE estructural} &
\textbf{Diferencia} \\
\midrule
2024 & $-3{,}46$ & $-3{,}94$ & $-0{,}48$ \\
\rowcolor{rowgray}
2025 & $-2{,}15$ & $-3{,}04$ & $-0{,}89$ \\
2026 & $-1{,}80$ & $-2{,}19$ & $-0{,}39$ \\
\rowcolor{rowgray}
2027 & $-1{,}40$ & $-1{,}74$ & $-0{,}34$ \\
2028 & $-1{,}00$ & $-1{,}28$ & $-0{,}28$ \\
\rowcolor{rowgray}
2029 & $-0{,}99$ & $-0{,}99$ & $\phantom{-}0{,}00$ \\
\bottomrule
\end{tabular}
\smallskip

\nota{Fuente: Cuadro~1 del MMM 2026--2029. El resultado estructural se expresa
en \% del PBI potencial; el convencional, en \% del PBI observado.}
\end{table}

Dos observaciones surgen de la Tabla~\ref{tab:estructural}:

\begin{enumerate}[leftmargin=1.5em]

    \item \emp{El resultado estructural es más negativo que el convencional
    en todos los años hasta 2028.}
    %
    Esto implica que el PBI observado se sitúa \emph{por encima} del PBI
    potencial --- la economía opera con una brecha de producto positiva ---,
    lo que infla transitoriamente la recaudación convencional. Dicho de otro
    modo, los ingresos reales del gobierno incluyen un componente cíclico
    favorable que no estará disponible cuando la economía retorne a su ritmo
    tendencial. La \emp{posición fiscal subyacente es más débil de lo que
    aparentan los números convencionales}.

    \item \emp{La brecha es máxima en 2025 ($-0{,}89$ pp) y desaparece en 2029.}
    %
    La convergencia de ambos indicadores en 2029 es consistente con el supuesto
    del MMM de que, al final del horizonte, el PBI observado habrá convergido
    a su nivel potencial (brecha del producto nula). En ese punto, todos los
    ingresos son de naturaleza estructural y la distinción entre los dos
    resultados desaparece.

\end{enumerate}

\begin{tcolorbox}[title={Implicancia de política fiscal}]
Si el gobierno diseñara su política de gasto tomando como referencia el
resultado económico \emph{convencional}, estaría incorporando al presupuesto
ingresos cíclicos que no son sostenibles. Cuando la brecha del producto se
cierre --- y con ella se reduzca la recaudación transitoria ---, el gobierno
se encontraría con un gasto estructuralmente más alto que sus ingresos
permanentes. El resultado estructural existe precisamente para evitar esa
trampa.
\end{tcolorbox}

% ------------------------------------------------------------
\subsection{Trayectoria del déficit: ¿qué dicen los números
            y qué los explica? --- inciso (e)}
% ------------------------------------------------------------

\subsubsection*{Descripción de la trayectoria histórica 2019--2025}

La serie histórica del BCRP (Cuadro~4 del repositorio) revela una trayectoria
con tres episodios bien delimitados en el período reciente:

\begin{table}[h!]
\centering
\caption{Resultado económico y primario del SPNF e ingresos del Gobierno
         Central (\% del PBI), 2019--2025}
\label{tab:historico}
\setlength{\tabcolsep}{8pt}
\renewcommand{\arraystretch}{1.25}
\begin{tabular}{lrrrrr}
\toprule
\rowcolor{headgray}
\textbf{Año} &
\textbf{RE} &
\textbf{RP} &
\textbf{Intereses} &
\textbf{Ing.\ Trib.\ GC} &
\textbf{IR (GC)} \\
\midrule
2019 & $-1{,}59$ & $-0{,}24$ & $1{,}35$ & $14{,}03$ & $5{,}58$ \\
\rowcolor{rowgray}
2020 & $-8{,}71$ & $-7{,}14$ & $1{,}57$ & $12{,}68$ & $5{,}20$ \\
2021 & $-2{,}48$ & $-1{,}00$ & $1{,}48$ & $15{,}68$ & $6{,}15$ \\
\rowcolor{rowgray}
2022 & $-1{,}68$ & $-0{,}13$ & $1{,}54$ & $16{,}56$ & $7{,}34$ \\
2023 & $-2{,}75$ & $-1{,}11$ & $1{,}64$ & $14{,}45$ & $6{,}17$ \\
\rowcolor{rowgray}
2024 & $-3{,}45$ & $-1{,}78$ & $1{,}67$ & $14{,}03$ & $5{,}93$ \\
2025* & $-2{,}18$ & $-0{,}59$ & $1{,}58$ & $14{,}42$ & $6{,}51$ \\
\bottomrule
\end{tabular}
\smallskip

\nota{* Cifras preliminares.
Fuente: Cuadro~4 del repositorio (BCRP). IR = Impuesto a la Renta del
Gobierno Central. Todas las variables en \% del PBI.}
\end{table}

Los tres episodios son:

\paragraph{Episodio 1 -- El choque de la pandemia (2020).}
El resultado económico se desplomó de $-1{,}59\%$ en 2019 a
$\mathbf{-8{,}71\%}$ del PBI en 2020, el peor registro de la historia
fiscal reciente. El deterioro fue casi íntegramente primario: el resultado
primario pasó de $-0{,}24\%$ a $-7{,}14\%$. Ello refleja la combinación de
una caída de ingresos tributarios (de 14{,}03\% a 12{,}68\% del PBI) y un
gasto extraordinario en transferencias, bonos y programas de apoyo al sector
privado (Reactiva Perú, entre otros). Los intereses apenas variaron, lo que
confirma que el choque fue de política fiscal discrecional, no de servicio
de deuda.

\paragraph{Episodio 2 -- La rápida recuperación (2021--2022).}
La combinación de rebote de actividad, alza de precios de materias primas
(especialmente cobre y zinc) y normalización del gasto extraordinario permitió
una recuperación fiscal veloz. El resultado económico mejoró de $-8{,}71\%$
a $-1{,}68\%$ en solo dos años. El IR del Gobierno Central subió de 5{,}20\%
a 7{,}34\% del PBI, impulsado por las ganancias extraordinarias de las empresas
mineras. En 2022, el resultado primario era ya casi equilibrado ($-0{,}13\%$).

\paragraph{Episodio 3 -- El nuevo deterioro (2023--2024).}
El resultado económico empeoró de $-1{,}68\%$ en 2022 a $-3{,}45\%$ en 2024,
revirtiéndose casi completamente la consolidación post-pandemia. Este es el
episodio más relevante para entender la posición fiscal que hereda el MMM
2026--2029 y se analiza en detalle a continuación.

\subsubsection*{¿Qué explica el deterioro 2022--2024?}

El deterioro fue principalmente un fenómeno de \emp{ingresos}, aunque agravado
por la rigidez del gasto corriente.

\bigskip
\noindent\textbf{Lado de los ingresos (Cuadro~6 del repositorio):}

Los ingresos tributarios del Gobierno Central cayeron de
$16{,}56\%$ del PBI en 2022 a $14{,}03\%$ en 2024, una contracción de
\emp{2{,}53 puntos porcentuales}. Tres factores la explican:

\begin{enumerate}[leftmargin=1.5em]
    \item \emp{Normalización del ciclo de materias primas.} El IR recaudado
          sobre Personas Jurídicas --- que recoge principalmente las ganancias
          mineras --- cayó de manera significativa al moderarse los precios
          del cobre. El IR total pasó de 7{,}34\% a 5{,}93\% del PBI entre
          2022 y 2024 ($-1{,}41$ pp), siendo la caída de personas jurídicas
          el factor dominante.
    \item \emp{Desaceleración de la actividad doméstica.} La recaudación del
          IGV, que mide el pulso del consumo interno, retrocedió de 9{,}27\%
          a 7{,}97\% del PBI ($-1{,}30$ pp), reflejando el menor dinamismo
          económico de 2023 (año marcado además por convulsión política y
          efectos rezagados del Fenómeno El Niño 2023).
    \item \emp{Aumento de devoluciones y retroceso de otros ingresos.}
          Las devoluciones tributarias y la caída de otros ingresos
          no tributarios contribuyeron al resultado adverso.
\end{enumerate}

\bigskip
\noindent\textbf{Lado del gasto (Cuadro~10 del repositorio):}

Los gastos no financieros del gobierno general se mantuvieron en torno a
$20{,}77\%$ del PBI en 2024, nivel superior al de años pre-pandemia
($\approx19\%$). La rigidez proviene principalmente de dos rubros:

\begin{enumerate}[leftmargin=1.5em]
    \item \emp{Remuneraciones (6{,}06\% del PBI en 2025 según C-10).} El gasto
          en planilla del sector público creció en términos nominales y
          resultó difícil de comprimir ante caídas de ingresos.
    \item \emp{Transferencias (2{,}72\% del PBI).} Las transferencias corrientes
          --- que incluyen pensiones, subsidios y programas sociales ---
          tienen un carácter cuasi-fijo en el corto plazo.
\end{enumerate}

La combinación de una caída pronunciada de ingresos con un gasto corriente
rígido es la explicación más directa del deterioro. El resultado primario
pasó de $-0{,}13\%$ (2022) a $-1{,}78\%$ (2024), y dado que los intereses
también aumentaron modestamente (de 1{,}54\% a 1{,}67\%), el resultado
económico acumuló una caída de casi $1{,}8$ puntos porcentuales.

\subsubsection*{Las proyecciones del MMM y sus condiciones de viabilidad}

El MMM proyecta una mejora sostenida hasta $-0{,}99\%$ del PBI en 2029.
Para que esa trayectoria se materialice deben cumplirse simultáneamente
las siguientes condiciones:

\begin{table}[h!]
\centering
\caption{Supuestos implícitos en la trayectoria fiscal proyectada
         (\% del PBI), 2024--2029}
\label{tab:supuestos}
\setlength{\tabcolsep}{9pt}
\renewcommand{\arraystretch}{1.25}
\begin{tabular}{lrrrrr}
\toprule
\rowcolor{headgray}
\textbf{Variable} & \textbf{2024} & \textbf{2026} & \textbf{2027}
                  & \textbf{2028} & \textbf{2029} \\
\midrule
Ingresos GG & $18{,}73$ & $19{,}15$ & $19{,}18$ & $19{,}19$ & $19{,}19$ \\
\rowcolor{rowgray}
Gastos no financieros GG & $20{,}77$ & $19{,}30$ & $18{,}91$ & $18{,}51$ & $18{,}55$ \\
Gastos corrientes & $14{,}68$ & $14{,}11$ & $13{,}70$ & $13{,}32$ & $13{,}16$ \\
\rowcolor{rowgray}
Gastos de capital & $6{,}09$ & $5{,}30$ & $5{,}22$ & $5{,}19$ & $5{,}39$ \\
Resultado primario & $-1{,}78$ & $-0{,}07$ & $+0{,}36$ & $+0{,}78$ & $+0{,}75$ \\
\rowcolor{rowgray}
Resultado económico & $-3{,}46$ & $-1{,}80$ & $-1{,}40$ & $-1{,}00$ & $-0{,}99$ \\
\bottomrule
\end{tabular}
\smallskip

\nota{Fuente: Cuadro~1 del MMM 2026--2029.}
\end{table}

\noindent La Tabla~\ref{tab:supuestos} permite identificar que la consolidación
descansa en dos pilares:

\begin{itemize}[leftmargin=1.5em]
    \item \emp{Recuperación de ingresos:} se proyecta que los ingresos del
          gobierno general se estabilicen en torno a 19{,}15--19{,}19\% del
          PBI, niveles similares a los de los años 2022--2023, bajo el supuesto
          de que la actividad económica se recupera y la recaudación del IR y
          del IGV se normaliza.
    \item \emp{Compresión del gasto corriente:} los gastos corrientes caerían
          $1{,}5$ puntos porcentuales entre 2024 y 2029, lo que implica que
          el crecimiento real del gasto corriente deberá ser
          \emph{sistemáticamente inferior} al crecimiento del PBI durante
          todo el período.
\end{itemize}

\medskip
\noindent Los principales \emp{riesgos} que podrían desviar esta trayectoria son:

\begin{enumerate}[leftmargin=1.5em]
    \item \emp{Nueva caída de precios de materias primas.} Dado que el IR de
          personas jurídicas es altamente sensible al ciclo minero, una
          corrección de precios del cobre o del zinc volvería a contraer la
          recaudación de manera abrupta.
    \item \emp{Rigidez del gasto corriente.} Comprimir las remuneraciones y
          las transferencias corrientes en términos reales requiere voluntad
          política sostenida y es difícil de mantener en contextos de
          inestabilidad institucional.
    \item \emp{Shocks externos o climáticos.} Un nuevo episodio de Fenómeno
          El Niño, una desaceleración global o tensiones geopolíticas podrían
          reducir tanto los ingresos como la capacidad de consolidar el gasto.
    \item \emp{Presiones de gasto inflexibles.} El aumento proyectado de
          intereses ($+0{,}07$ pp entre 2024 y 2029) es modesto, pero una
          depreciación cambiaria significativa elevaría el costo de la deuda
          externa y añadiría presión sobre el resultado económico.
\end{enumerate}


% ============================================================
\section{Las tres etapas de la política fiscal peruana\\
         \large Respuesta a Tito Marx}
% ============================================================

\noindent La siguiente caracterización se apoya en los dos gráficos del archivo
\texttt{the\_fiscal\_history\_of\_PER\_in\_two\_graphs.pdf} elaborado por el
Ministerio de Economía y Finanzas y el Banco Central de Reserva del Perú.
Las figuras reportan, respectivamente, el resultado económico y primario del
SPNF (en \% del PBI) y el saldo de deuda pública (en \% del PBI) para el
período 1970--2017.

% ------------------------------------------------------------
\subsection{Primera etapa: la crisis fiscal de los años 1980}
% ------------------------------------------------------------

El primer gráfico muestra que los años 1980 constituyen la \emp{peor década
fiscal en la historia contemporánea del Perú}. El resultado económico se
deterioró de manera progresiva y alcanzó un piso cercano a \emp{$-12\%$ del
PBI} hacia 1983--1985, acompañado de un resultado primario también
profundamente negativo (alrededor de $-4\%$ a $-5\%$ del PBI). La deuda
pública, que partía de aproximadamente 13\% del PBI en 1970, escaló
vertiginosamente hasta un pico de \emp{96{,}2\% del PBI} al filo de
1988--1990, según el segundo gráfico.

El gráfico identifica dos eventos que condensan la historia de la década:

\begin{enumerate}[leftmargin=1.5em]
    \item \emp{El Niño de 1982--83.} El fenómeno climático golpeó la producción
          agropesquera y energética, redujo los ingresos fiscales y amplió el
          déficit en un contexto en que las cuentas externas ya eran frágiles
          por el inicio de la crisis de deuda latinoamericana.

    \item \emp{Repudiación de la deuda, controles de precios e hiperinflación
          (segunda mitad de la década).} Este es el evento central. Bajo el
          gobierno de Alan García (1985--1990), el Perú adoptó una combinación
          de políticas que resultó catastrófica: se declaró que el servicio de
          la deuda externa no superaría el 10\% de las exportaciones; se
          implantaron controles de precios para contener la inflación; y el
          déficit fiscal se financió en creciente medida con emisión monetaria.
          El resultado fue una hiperinflación que superó el 7\,000\% anual
          en 1990 y la explosión de la deuda pública hasta 96{,}2\% del PBI.
\end{enumerate}

\begin{tcolorbox}[title={Un apunte para Tito Marx}]
La hiperinflación es el impuesto más regresivo que existe: destruye el poder
adquisitivo de los trabajadores y de quienes no poseen activos reales
con que protegerse. El daño económico más severo de los años 1980 recayó,
precisamente, sobre los sectores populares. La irresponsabilidad fiscal
---independientemente de la orientación política que la impulse--- termina
empobreciendo a los más pobres.
\end{tcolorbox}

% ------------------------------------------------------------
\subsection{Segunda etapa: estabilización y construcción
            institucional en los años 1990}
% ------------------------------------------------------------

El primer gráfico muestra una recuperación abrupta del resultado primario
alrededor de 1990--1991, etiquetada como \emp{programa de estabilización}.
El resultado primario pasó de terreno profundamente negativo a positivo
($+2\%$ del PBI) en el espacio de dos o tres años. Se trata del ajuste fiscal
de mayor velocidad en toda la serie histórica.

Sin embargo, el resultado económico \emph{global} permaneció negativo a lo
largo de toda la década. La explicación es precisa: el déficit remanente era
explicado casi íntegramente por los intereses devengados sobre la deuda
acumulada en los años 1980, no por nuevo gasto primario. El gobierno había
logrado equilibrar su gasto corriente, pero todavía cargaba el peso del
endeudamiento heredado.

El gráfico registra tres hitos institucionales que definen esta etapa:

\begin{enumerate}[leftmargin=1.5em]
    \item \emp{Programa de estabilización (1990--1991).} El denominado
          ``Fujishock'' combinó liberalización de precios, unificación
          cambiaria y una fuerte contracción fiscal. La inflación cedió
          rápidamente.
    \item \emp{Nueva constitución y reformas (1993).} La Constitución de 1993
          introdujo principios de equilibrio presupuestal y sentó las bases
          legales para una nueva arquitectura fiscal. Junto con las
          privatizaciones y la reforma tributaria (creación de la SUNAT),
          permitió reconstruir la base de ingresos del Estado.
    \item \emp{Ley de Responsabilidad y Transparencia Fiscal e introducción
          de reglas fiscales ($\approx$1999).} El Perú formalizó un marco de
          reglas numéricas sobre el déficit y el gasto, convirtiéndose en uno
          de los primeros países de la región en adoptar esta herramienta.
\end{enumerate}

El segundo gráfico muestra el resultado de esta etapa sobre la deuda: desde
el pico de 96{,}2\% del PBI, el stock de deuda descendió de forma sostenida
hasta situarse en \emp{48{,}7\% del PBI} a comienzos de los años 2000.
Notablemente, el efecto El Niño de 1997--98 provocó una interrupción temporal
visible en ambas gráficas, pero no comprometió la tendencia descendente.

% ------------------------------------------------------------
\subsection{Tercera etapa: consolidación, ahorro del ciclo y
            desgaste gradual en las décadas de 2000 y 2010}
% ------------------------------------------------------------

Esta etapa exhibe dos momentos fiscalmente muy distintos, perfectamente
legibles en los gráficos.

\paragraph{El boom de materias primas y el ahorro del superávit (2004--2008).}
El primer gráfico muestra cómo el resultado primario ascendió hasta
aproximadamente \emp{$+5\%$ del PBI} hacia 2007, generando superávits
globales por primera vez en décadas. La anotación \emp{``Saved surplus:
Fiscal Stabilization Fund''} es la clave interpretativa: a diferencia de
episodios anteriores, el Estado no utilizó el ingreso extraordinario para
expandir el gasto de forma permanente, sino que lo esterilizó parcialmente
en el Fondo de Estabilización Fiscal. Esta fue la decisión de política fiscal
más importante de la etapa y explica en gran medida la resiliencia que el
Perú mostró ante la crisis financiera global de 2008--2009.

El segundo gráfico recoge el efecto sobre el stock de deuda: la combinación
de superávits primarios y crecimiento elevado redujo la deuda pública de
48{,}7\% a niveles cercanos al \emp{19--20\% del PBI} hacia 2012--2013,
para luego cerrar el período en \emp{25{,}8\% del PBI}.

\paragraph{Fortalecimiento institucional y regreso a la moderación (2009--2017).}
El gráfico registra dos nuevos hitos en esta sub-etapa:

\begin{itemize}[leftmargin=1.5em]
    \item \emp{Reforma de pensiones} (alrededor de 2000--2001): restructuración
          del sistema previsional que generó costos fiscales transitorios pero
          mejoró la sostenibilidad de largo plazo.
    \item \emp{Fortalecimiento del marco fiscal: balance estructural y
          consejo fiscal} ($\approx$2012--2013): el Perú adoptó la regla de
          balance estructural --- el mismo instrumento cuya lógica se analiza
          en el inciso (d) --- y creó el Consejo Fiscal como instancia de
          evaluación independiente.
    \item \emp{Transparencia y retorno a metas nominales} ($\approx$2016):
          tras la experiencia con el balance estructural, el marco volvió a
          complementarse con anclas nominales explícitas sobre el déficit y
          la deuda.
\end{itemize}

Cuando el boom de materias primas cedió (a partir de 2012), el resultado
fiscal comenzó a deteriorarse nuevamente de forma gradual. El resultado
económico pasó de superávit a déficits moderados de $-1\%$ a $-2\%$ del PBI
hacia el final del período. Esto reveló que parte del gasto expandido durante
el auge se había vuelto estructural: al caer los ingresos cíclicos, el
presupuesto no pudo ajustarse con igual rapidez.

\begin{tcolorbox}[title={Síntesis para Tito Marx}]
El arco completo de los dos gráficos transmite una lección que trasciende las
divisiones ideológicas: las tres etapas demuestran que la disciplina fiscal
no es un fin en sí mismo, sino la condición previa para que el Estado pueda
cumplir sus objetivos sociales de manera sostenible. La etapa de mayor
expansión del gasto sin respaldo fiscal (1980s) produjo hiperinflación y
96{,}2\% de deuda; la etapa de mayor disciplina y ahorro (2000s) permitió
reducir la deuda a menos de un cuarto de ese pico y financiar programas
sociales sin crear vulnerabilidades. \emp{Gastar bien y gastar de forma
sostenible no son objetivos en tensión: son el mismo objetivo}.
\end{tcolorbox}

% ============================================================
\section*{Fuentes de datos}
% ============================================================

\begin{itemize}[leftmargin=1.5em]
    \item Ministerio de Economía y Finanzas del Perú (MEF). \textit{Marco
          Macroeconómico Multianual 2026--2029}. Lima, 2025. Cuadros n.\textordmasculine~6,
          8, 9, 12, 13, 14, 15, 18 y 19.
    \item Banco Central de Reserva del Perú (BCRP). \textit{BCRPData ---
          Base de Datos de Estadísticas}. Series anuales: resultado económico
          (PM05780FA), resultado primario (PM05776FA), intereses (PM10083FA),
          ingresos tributarios del gobierno central (PM10103FA), IGV
          (PM10108FA), IR (PM10104FA), deuda pública total (PM10188FA), entre
          otras. Período 2000--2025.
    \item MEF / BCRP. \textit{The fiscal history of Peru in two graphs}.
          Archivo \texttt{the\_fiscal\_history\_of\_PER\_in\_two\_graphs.pdf}.
    \item Segura, A. (2015). ``Some Thoughts on Fiscal Policy and the
          Unfinished Agenda.'' En \textit{Peru: Staying the Course of Economic
          Success}, cap.~24. FMI.
    \item Rossini, R. y Santos, A. (2015). ``Peru's Recent Economic History.''
          En \textit{Peru: Staying the Course of Economic
          Success}, cap.~2. FMI.
\end{itemize}

\end{document}
